% META: {"id": "1SPE-TRIGO-CO-022", "chapitre": "1SPE-TRIGONOMETRIE", "exercice_ref": "1SPE-TRIGO-EX-022", "type_objet": "corrige", "capacites": ["1SPE-TRIGONOMETRIE-C1"], "capacites_codes": ["C1"], "methodes": ["M1"], "parcours": 1, "competences": ["calculer"], "duree_min": 5, "mode_creation": "ex_nihilo", "sources_inspiration": [], "parametres_sympy": {}, "status": "approved"}
\begin{corrige}{1SPE-TRIGO-EX-022}
$\dfrac{17\pi}{3} = 6\pi - \dfrac{\pi}{3}$. La mesure principale est $-\dfrac{\pi}{3}\text{ rad}$.
\end{corrige}

% BEGIN-VERIFY
% from sympy import pi, Rational, simplify, cos, sin
%
% angle = Rational(17, 3) * pi
% principale = -pi/3
%
% # La mesure principale appartient a ]-pi ; pi] et differe de l'angle d'un
% # multiple entier de 2pi.
% assert -pi < principale <= pi
% ecart = simplify((angle - principale) / (2*pi))
% assert ecart == 3 and ecart.is_integer
%
% # Deux angles de meme mesure principale ont memes cosinus et sinus.
% assert simplify(cos(angle) - cos(principale)) == 0
% assert simplify(sin(angle) - sin(principale)) == 0
%
% # La decomposition annoncee, 17pi/3 = 6pi - pi/3, est exacte.
% assert simplify(angle - (6*pi - pi/3)) == 0
% END-VERIFY
